\section{Related Work}
\label{sec_Related_Work}

\noindent \textbf{3D visual grounding.}
Significant progress has been made in 3D visual grounding thanks to advancements in deep learning and computer vision, as well as the availability of grounded datasets \citep{achlioptas2020referit3d, chen2020scanrefer, abdelreheem2022scanents3d, yuan2022toward}. 
One approach involves using graph-based models \citep{achlioptas2020referit3d, feng2021free, yuan2021instancerefer, huang2021text} to represent the scene as nodes and edges, while attention mechanisms help to focus on relevant parts of the scene.
Another approach \citep{roh2022languagerefer} is to convert the visual input into language tokens using a classification head. These tokens can then be combined with the input utterance and fed into a transformer architecture to learn the relationships between input sequence elements.
Moreover, recent work \citep{bakr2022look, yang2021sat} explores distilling knowledge from 2D to 3D in a multi-view setup. 
However, none of the existing works models the explicit reasoning process behind the prediction of the target object.

\noindent \textbf{Chain of thoughts.}
The chain-of-thought concept has been used in many different machine learning applications, including natural language processing \citep{wei2022chain, chowdhery2022palm, lyu2023faithful, wang2022self, zhang2023multimodal, madaan2022text}, and robotics \citep{jia2023chain, yang2022chain}. 
In the context of 3D Visual grounding, developing a chain-of-thought approach provides a natural way to explicitly model the grounding reasoning process, which to the best of our knowledge has not been explored.
An extended version is discussed in the Appendix \ref{appendix_related_work}.